\documentclass[french]{article}
\usepackage[utf8]{inputenc}
\usepackage[T1]{fontenc}
\usepackage{lmodern}
\usepackage{geometry}
\usepackage{graphicx}
\usepackage{babel}
\usepackage{amsmath}
\usepackage{amsthm}
\usepackage{amssymb}
\usepackage{hyperref}
\usepackage{stmaryrd}
\usepackage{enumitem}
\usepackage{tcolorbox}
\usepackage{dsfont}
\usepackage{multirow}
\usepackage{etoc}
\usepackage{titlesec}
\usepackage{esint}
\usepackage{cancel}
\usepackage{mathdots}

\setlist[itemize]{label=$\bullet$}


\renewcommand{\P}{\mathbb{P}}
\newcommand{\N}{\mathbb{N}}
\newcommand{\Z}{\mathbb{Z}}
\newcommand{\Q}{\mathbb{Q}}
\newcommand{\R}{\mathbb{R}}
\newcommand{\C}{\mathbb{C}}
\newcommand{\K}{\mathbb{K}}
\renewcommand{\L}{\mathbb{L}}
\newcommand{\U}{\mathbb{U}}
\newcommand{\st}{^*}
\newcommand{\tq}{\middle|}
\newcommand{\inv}{^{-1}}
\newcommand{\E}{\mathbb{E}}
\newcommand{\F}{\mathbb{F}}
\newcommand{\V}{\mathbb{V}}
\renewcommand{\ker}{\text{Ker}}
\newcommand{\im}{\text{Im}}
\newcommand{\card}{\text{Card}}
\newcommand{\Tr}{\text{Tr}}

\title{TD Euclidiens\\ Exercice 11}
\date{}
\begin{document}
\maketitle

\section{Énoncé}
\begin{enumerate}
	\item Montrer que toute matrice inversible $M$ peut s'écrire $M=OT$ avec $O$ orthogonale et $T$ triangulaire supérieure.
	\item Montrer l'unicité d'un tel couple si on impose que les coefficients diagonaux de $T$ soient tous strictement positifs.
	\item Mêmes questions pour la décomposition $TO$.
	\item Application : inégalité de Hadamard. Montrer : $$\lvert\det(M)\rvert\leqslant\prod_{j=1}^{n}\lVert C_j\rVert$$ où les $C_j$ sont les colonnes de $M$ et où $\lVert\cdot\rVert$ désigne la norme euclidienne canonique sur $\R^n$.
	\item Que reste-t-il de ces résultats pour les matrices non inversibles ?
\end{enumerate}
\section{Solution}
\begin{enumerate}
	\item Appliquer le processus de Gram-Schmidt aux colonnes de $M$ et effectuer un changement de base.
	\item Utiliser le fait que les matrices triangulaires et orthogonales sont des diagonales de $\pm 1$ et conclure avec les strictes positivités.
	\item Idem, mais à l'envers (ou transposer).
	\item Reprendre les étapes de la décompositions $OT$ et utiliser l'inégalité de Cauchy-Schwarz.
	\item L'existence des décompositions reste vraie par compacité de $\mathcal{O}_n(\R)$ mais pas l'unicité. L'inégalité de Hadamard est trivialement vraie.
\end{enumerate}

\end{document}
